\section{Threats to validity and what remains}
\label{sec:threats}

\paragraph{We do not report a corrected ranking.}
The corrected infrastructure exists and every ecosystem has been verified on the
accelerator, but the replicated campaign has not been executed at the time of
writing. We consider this the honest position: a ranking produced before the
defects of Sections~\ref{sec:instrument} and~\ref{sec:implementation} are fixed
would inherit them, and a ranking produced after them is a different experiment
that deserves its own reporting.

\paragraph{Residual confounds we could not remove.}
Two survive by construction rather than by oversight. DL4J 1.0.0-M2.1 links
against the CUDA 11.6 runtime and is the last release of its API, so the Java
ecosystem cannot be brought onto the CUDA 12 toolchain the other six use; the
stack now carries its own redistributables, which makes it reproducible but does
not remove the version difference. And the R binding cannot switch a script
module between training and evaluation, so it cannot share the traced model.
Both are properties of the ecosystems, and arguably part of what one is measuring
when one measures an ecosystem --- but they are not the binding overhead the
study set out to isolate.

\paragraph{The workload.}
The campaign runs the accelerator at \SIrange{24}{56}{\percent} of its power
limit, so it measures data loading and dispatch more than deep learning. A
replicated campaign should add at least one configuration that saturates the
accelerator, and report utilisation alongside energy so a reader can see which
regime a result belongs to.

\paragraph{The instrument itself.}
We now read hardware counters directly ---
\code{nvmlDeviceGetTotalEnergyConsumption} for the accelerator and Intel RAPL for
the CPU package --- and run CodeCarbon over the same window as a cross-check.

It is worth being precise about what was wrong with the original instrument,
because it was not accuracy. Measured on the replication machine with both
instruments bracketing the same window, CodeCarbon 3.3.0 agrees with the counters
to within \SI{1}{\percent} for block durations from \SI{2}{\second} to
\SI{40}{\second}, for GPU and CPU alike. The failure mode is that it
\emph{substitutes a model when a counter is unavailable, and says so only in a
log line}: where RAPL cannot be read the CPU term becomes a constant, the RAM term
is derived from installed memory rather than from use, and the model changed
between major versions. An instrument that refused to run would have been
preferable to one that quietly estimated.

This matters for how a number should be reported. NVML plus RAPL is a
\emph{chip-level} boundary: it excludes the power supply, fans and drives that a
wall meter would see. Published validations place software estimators within tens
of percent of a wall meter, so a chip-boundary figure must not be presented as a
whole-system one. The figure the campaign under audit reports sits between the two
boundaries and corresponds to neither --- for four ecosystems roughly two thirds
of it was a model of RAM power derived from how much memory the machine happened
to have installed.

We had no wall meter, and say so rather than implying a boundary we did not
measure.

\paragraph{Generalisation.}
Everything here concerns one campaign, one hardware platform and one model class.
The defect \emph{classes} of Section~\ref{sec:catalogue} are, we believe, general
to studies of this shape; the effect sizes are not.

\paragraph{Our own errors.}
Two of the defects in Section~\ref{sec:catalogue} we introduced ourselves while
repairing the study, and caught only because the accuracy check of
Table~\ref{tab:convergence} existed: switching the Rust stack to the shared
TorchScript module made it evaluate in training mode, costing 15 accuracy points,
because \code{TrainableCModule} ignores the training flag passed to
\code{forward\_t}. We measured contaminated blocks ourselves by running package
downloads while the campaign was tracking whole-machine energy. We report this
because it is the argument of the paper in miniature: the defects that matter are
the ones that produce plausible numbers, and the only reliable defence is a check
that fails.
